\textit{Pour chacune des affirmations suivantes, indiquer si elle est \textbf{vraie} ou \textbf{fausse}, en \textbf{justifiant} la réponse. Un réponse non justifiée n'est pas prise en compte.}

\begin{center}
	\textit{Les trois questions sont indépendantes.}
\end{center}

\begin{enumerate}
	\item On considère un \textbf{réel} $u$. On considère sur $\mathbb{R}$ l'équation

	\begin{center}
		(\textit{E})\hspace{10mm}\mbox{$x^2+x-u^2=0$}
		\end{center}

		\textbf{Affirmation} : Quelle que soit la valeur du réel $u$, l'équation (\textit{E}) possède deux solutions réelles distinctes.


		\textcolor{red}{\textsc{Correction}}


		Le discriminant de ce polynome est :\quad \mbox{$1+4u^2$}.
		il est donc strictement positif pour toute valeur de $u$, et l'équation (\textit{E}) a toujours deux solutions,

		\hfill \framebox[1.2\width]{affirmation vraie.}

		\item On considère la suite \mbox{$(u_n)$} définie pour tout entier naturel $n$ par

		\begin{center}
			\mbox{$u_n=2^{-n}$}
		\end{center}

		\textbf{Affirmation} : La suite \mbox{$(u_n)$} est une suite géométrique de raison \mbox{$\dfrac12$}.


		\textcolor{red}{\textsc{Correction}}


		\mbox{$\forall n\in\N, \quad u_n=2^{-n}=\dfrac{1}{2^n}=\dfrac{1^n}{2^n}=\left(\dfrac{1}{2}\right) ^n$} :\quad \mbox{$(u_n)$} est bien une suite géométrique de raison \mbox{$\dfrac12$}.

		Son premier terme est \quad $u_0 = 2^{-0} = 2^0 = 1$,

		\hfill \framebox[1.2\width]{affirmation vraie.}

		\item On considère la fonction $f$ définie sur $\mathbb{R}$ par \mbox{$f(x)=\e^x-1$}.

		On note \mbox{$\mathcal{C}$} la courbe représentative de la fonction $f$ dans un repère.

		On note \mbox{$\mathcal{T}$} la tangente à la courbe \mbox{$\mathcal{C}$} au point d'abscisse $0$.

		On note A le point de coordonnées \mbox{$(3;3)$.}

		\vspace{5mm}

		\textbf{Affirmation} : Le point A appartient à la tangente \mbox{$\mathcal{T}$}.


		\textcolor{red}{\textsc{Correction}}

		La formule pour avoir l'équation réduite de la tangente à la courbe représentative d'une fonction $f$ au point d'abscisse $a$ est :\quad $y = f'(a)\times (x - a) + f(a)$.

		On va l'appliquer ici avec $a = 0$.

		La dérivée de $f$ est \quad \mbox{$f'(x)=\e^{x}$}.

		Au point d'abscisse $0$, le nombre dérivé est :\quad $f'(0) = \e^{0} = 1$

		et\quad  $f(0)=0$.

		Donc l'équation de \mbox{$\mathcal{T}$} est \quad \mbox{$y=x$}.

		le point A appartient à la tangente, car ses coordonnées en vérifient l'équation.

		\hfill \framebox[1.2\width]{affirmation vraie.}
\end{enumerate}
